\documentclass[aspectratio=169,10pt]{beamer}

% University of Tehran Color Theme & Styling
\usetheme{Madrid}
\usecolortheme{whale}

\definecolor{UTBlue}{RGB}{12, 44, 110}       % University of Tehran Deep Blue
\definecolor{UTGold}{RGB}{200, 150, 40}      % UT Accent Gold
\definecolor{LightBlue}{RGB}{240, 245, 255}
\definecolor{TealDark}{RGB}{0, 100, 100}

\setbeamercolor{palette primary}{bg=UTBlue,fg=white}
\setbeamercolor{palette secondary}{bg=UTBlue!85!black,fg=white}
\setbeamercolor{palette tertiary}{bg=UTBlue!70!black,fg=white}
\setbeamercolor{titlelike}{parent=palette primary}
\setbeamercolor{structure}{fg=UTBlue}

\setbeamercolor{block title}{bg=UTBlue,fg=white}
\setbeamercolor{block body}{bg=LightBlue}
\setbeamercolor{block title example}{bg=TealDark,fg=white}
\setbeamercolor{block body example}{bg=teal!5!white}
\setbeamercolor{block title alerted}{bg=red!80!black,fg=white}
\setbeamercolor{block body alerted}{bg=red!5!white}

\setbeamertemplate{navigation symbols}{}
\setbeamertemplate{headline}{}

\setbeamertemplate{footline}{
  \leavevmode%
  \hbox{%
  \begin{beamercolorbox}[wd=.333333\paperwidth,ht=2.25ex,dp=1ex,center]{author in head/foot}%
    \usebeamerfont{author in head/foot}\insertshortauthor
  \end{beamercolorbox}%
  \begin{beamercolorbox}[wd=.333333\paperwidth,ht=2.25ex,dp=1ex,center]{title in head/foot}%
    \usebeamerfont{title in head/foot}\insertshorttitle
  \end{beamercolorbox}%
  \begin{beamercolorbox}[wd=.333333\paperwidth,ht=2.25ex,dp=1ex,center]{date in head/foot}%
    \insertframenumber{} / \inserttotalframenumber
  \end{beamercolorbox}}%
  \vskip0pt%
}

\usepackage{amsmath,amssymb,amsthm}
\usepackage{booktabs}
\usepackage{multirow}
\usepackage{graphicx}
\usepackage{tikz}
\usetikzlibrary{arrows.meta,positioning,shapes,calc}

% Persian Package Configuration
\usepackage{xepersian}
\settextfont[Path=fonts/,BoldFont={X Nazanin Bold.ttf},ItalicFont={X Nazanin Italic.ttf}]{X Nazanin.ttf}
\setlatintextfont[Scale=0.85]{Times New Roman}

\newcommand{\code}[1]{\lr{\texttt{#1}}}

\title[طراحی، ارزیابی و فشرده‌سازی توجه خطی و SVD]{طراحی، ارزیابی تجربی و فشرده‌سازی مکانیزم‌های توجه خطی و تجزیه مقادیر منفرد (SVD) در مدل‌های زبانی بزرگ}
\subtitle{\lr{A Hardware-Aware Framework for SVD Linear Attention and Transformer Model Compression}}
\author[محمد طاها مجلسی]{محمد طاها مجلسی \\ \small استاد راهنما: جناب آقای دکتر احمد کلهر}
\institute[دانشگاه تهران]{دانشکده مهندسی برق و کامپیوتر \\ پردیس دانشکده‌های فنی --- دانشگاه تهران}
\date{شهریور ۱۴۰۵}

\begin{document}

% --- Slide 1: Title Slide ---
\begin{frame}[plain]
  \begin{center}
    \includegraphics[height=1.25cm]{logo.png} \hspace{1.5cm}
    \includegraphics[height=1.25cm]{eng-logo.png} \\[0.3cm]
    {\large \textbf{دانشگاه تهران --- پردیس دانشکده‌های فنی}}\\[0.1cm]
    {\small دانشکده مهندسی برق و کامپیوتر}\\[0.3cm]
    {\large \color{UTBlue} \textbf{جلسه دفاع از پایان‌نامه مهندسی کامپیوتر}}\\[0.25cm]
    {\Large \textbf{طراحی، ارزیابی تجربی و فشرده‌سازی مکانیزم‌های توجه خطی و تجزیه مقادیر منفرد (SVD) در مدل‌های زبانی بزرگ}}\\[0.15cm]
    {\footnotesize \lr{A Hardware-Aware Framework for SVD Linear Attention and Transformer Compression}}\\[0.35cm]
    \begin{tabular}{rl}
      \textbf{نگارنده:} & محمد طاها مجلسی (\lr{810101504}) \\
      \textbf{استاد راهنما:} & جناب آقای دکتر احمد کلهر
    \end{tabular}
  \end{center}
\end{frame}

% --- Slide 2: Outline ---
\begin{frame}{سرفصل مطالب و ساختار ارائه}
  \begin{itemize}\setlength{\itemsep}{0.28cm}
    \item \textbf{انگیزه پژوهش:} دیوار حافظه ترنسفورمرها و پیچیدگی درجه دوم $\mathcal{O}(N^2)$
    \item \textbf{مبانی نظری و ریاضی:} سفیدسازی فعال‌سازی چولسکی و اکرت-یانگ-میرسکی
    \item \textbf{بسته‌شدن ۵ شکاف علمی مرز دانش (SOTA Gaps):} قضایای تحلیلی و تایید تجربی
    \item \textbf{پروتکل مهندسی تزریق جراحی ۴ مرحله‌ای} و هم‌ترازی گام صفر
    \item \textbf{بسته ۶ ستونی تقطیر دانش (6-Pillar KD Suite)} و اثبات انقباض خطی
    \item \textbf{ارزیابی پارتو در مقیاس وسیع:} ۲۰۳ ماژول (۱۳۱ توجه، ۷۲ فشرده‌ساز)
    \item \textbf{شتاب‌دهی سخت‌افزاری تریتون:} ۲۶ کرنل سفارشی \lr{Triton JIT} روی \lr{RTX 4070 SM 8.9}
    \item \textbf{دستاوردها، جمع‌بندی و نتیجه‌گیری}
  \end{itemize}
\end{frame}

% --- Slide 3: Motivation & Memory Wall ---
\begin{frame}{انگیزه پژوهش: دیوار حافظه و پیچیدگی $\mathcal{O}(N^2)$}
  \begin{columns}
    \begin{column}{0.55\textwidth}
      \begin{block}{چالش اصلی خودتوجهی متراکم (Softmax Attention)}
        محاسبه ماتریس ضرب داخلی پرسمان و کلید:
        \[
          A = \operatorname{Softmax}\left(\frac{Q K^\top}{\sqrt{d_k}}\right) V
        \]
        \begin{itemize}\footnotesize
          \item پیچیدگی زمانی و حافظه‌ای $\mathcal{O}(N^2)$ نسبت به طول توالی $N$.
          \item در توالی‌های فرابزرگ ($N > 32\text{K}$)، حافظه مادی‌سازی توجه از حجم فیزیکی \lr{VRAM} فراتر می‌رود.
          \item بحران بافر \lr{KV-Cache} در فاز استنتاج خودبازگشتی: مقیاس $\mathcal{O}(L \cdot b \cdot N \cdot d)$.
        \end{itemize}
      \end{block}
    \end{column}
    \begin{column}{0.45\textwidth}
      \begin{alertblock}{گلوگاه پهنای باند حافظه (Memory Wall)}
        \begin{itemize}\footnotesize
          \item سرعت تبادل داده میان \lr{HBM/DRAM} و ثبات‌های \lr{SRAM} گلوگاه اصلی اجرای مدل‌هاست.
          \item مدل‌های زبانی به شدت در ناحیه «محدود به حافظه» (\lr{Memory-Bound}) کار می‌کنند.
          \item راهکارهای سه‌گانه این رساله:
            \begin{enumerate}\scriptsize
              \item الگوریتم‌های توجه خطی $\mathcal{O}(N)$
              \item فشرده‌سازی مقادیر منفرد \lr{SVD} آگاه از فعال‌سازی
              \item ادغام کرنل‌های سخت‌افزاری \lr{Kernel Fusion}
            \end{enumerate}
        \end{itemize}
      \end{alertblock}
    \end{column}
  \end{columns}
\end{frame}

% --- Slide 4: Theoretical Foundations - Cholesky & SVD ---
\begin{frame}{مبانی نظری: سفیدسازی فعال‌سازی با تجزیه چولسکی}
  \begin{block}{فرمول‌بندی بهینه‌سازی برش مقادیر منفرد آگاه از فعال‌سازی}
    \small
    هدف، یافتن ماتریس کم‌مرتبه $\mathbf{W}_r$ با رتبه $r$ است که خطای فعال‌سازی خروجی را کمینه کند:
    \vspace{-0.15cm}
    \[
      \min_{\operatorname{rank}(\mathbf{W}_r) \le r} \mathbb{E}_{X \sim \mathcal{D}}\left[ \| \mathbf{W} X - \mathbf{W}_r X \|_F^2 \right] = \min_{\operatorname{rank}(\mathbf{W}_r) \le r} \| (\mathbf{W} - \mathbf{W}_r) \mathbf{H}_X^{1/2} \|_F^2
    \]
    \vspace{-0.15cm}
    که در آن $\mathbf{H}_X = \mathbb{E}[X X^\top]$ ماتریس کوواریانس ورودی‌های کالیبراسیون است.
  \end{block}

  \begin{exampleblock}{قضیه بهینگی SVD-LLM v2 با تجزیه چولسکی (Eckart-Young-Mirsky)}
    \small
    با تجزیه چولسکی کوواریانس $\mathbf{H}_X = \mathbf{L} \mathbf{L}^\top$ ($\mathbf{L}$ مثلثی پایینی و معکوس‌پذیر):
    \vspace{-0.15cm}
    \[
      \| (\mathbf{W} - \mathbf{W}_r) \mathbf{H}_X^{1/2} \|_F^2 = \| (\mathbf{W} - \mathbf{W}_r) \mathbf{L} \|_F^2 = \| \widetilde{\mathbf{W}} - \widetilde{\mathbf{W}}_r \|_F^2, \quad \widetilde{\mathbf{W}} = \mathbf{W} \mathbf{L}
    \]
    \vspace{-0.15cm}
    بنابر قضیه اکرت-یانگ-میرسکی، بهترین تقریب رتبه-$r$ ماتریس سفیدشده $\widetilde{\mathbf{W}}$ حاصل برش مقادیر منفرد آن است:
    \vspace{-0.15cm}
    \[
      \mathbf{W}_r^* = \widetilde{\mathbf{W}}_r \mathbf{L}^{-1} = \left(\sum_{i=1}^r \sigma_i u_i v_i^\top\right) \mathbf{L}^{-1}
    \]
  \end{exampleblock}
\end{frame}

% --- Slide 5: SVD-LLM v2 vs Naive Truncation Plot ---
\begin{frame}{اعتبارسنجی قضیه بهینگی SVD-LLM v2 در برابر برش مستقیم}
  \begin{columns}
    \begin{column}{0.55\textwidth}
      \includegraphics[width=\textwidth]{svd_llm_v2_truncation_optimality.png}
    \end{column}
    \begin{column}{0.45\textwidth}
      \begin{block}{تحلیل نمودار تجربی اعتبارسنجی}
        \begin{itemize}\footnotesize
          \item \textbf{برش سنتی بدون داده (Naive SVD):} انحراف شدید فروبنیوس به دلیل نادیده‌گرفتن همبستگی ورودی‌ها.
          \item \textbf{روش پیشنهادی SVD-LLM v2:} با سفیدسازی چولسکی، انحراف فعال‌سازی‌ها در تمام سطوح فشرده‌سازی تا ۱۰۰٪ مهار شده و خطای فروبنیوس به صفر میل می‌کند.
          \item تطابق دقیق نتایج عملی با کران تئوریک قضیه اکرت-یانگ-میرسکی.
        \end{itemize}
      \end{block}
    \end{column}
  \end{columns}
\end{frame}

% --- Slide 6: 5 Closed SOTA Gaps Table ---
\begin{frame}{بسته‌شدن قطعی ۵ شکاف علمی در مرز دانش (SOTA Gaps)}
  \begin{table}
    \centering
    \footnotesize
    \begin{tabular}{clllc}
      \toprule
      \textbf{شناسه} & \textbf{ماژول و کِرنل} & \textbf{موضوع علمی و قضیه تحلیلی} & \textbf{شاخص اعتبارسنجی تجربی} & \textbf{وضعیت} \\
      \midrule
      \textbf{Gap 1} & \lr{SVD-LLM v2} & بهینگی برش اکرت-یانگ-میرسکی چولسکی & کاهش خطای فروبنیوس: $28.68 \to 0.00$ & \textbf{\lr{PASSED}} \\
      \textbf{Gap 2} & \lr{ASVD Linear} & استریمینگ کم‌مرتبه بافر KV-Cache & تطابق عددی دقیق، کاهش حافظه $8\times$ & \textbf{\lr{PASSED}} \\
      \textbf{Gap 3} & \lr{FP8 Quantizer} & کران تحلیلی خطای مانتیس E4M3 & خطای نسبی $0.0264$ (زیر سقف $5\%$) & \textbf{\lr{PASSED}} \\
      \textbf{Gap 4} & \lr{DSA Attention} & مسیریابی تنک بلوکی پویا & هم‌ترازی گرادیان $>95\%$، برابری ابعاد & \textbf{\lr{PASSED}} \\
      \textbf{Gap 5} & \lr{REAP-MoE} & هرس نااریب خبرگان در توزیع بلند & هرس ۴ به ۲ خبره، بازیابی $100\%$ آزمون NIAH & \textbf{\lr{PASSED}} \\
      \bottomrule
    \end{tabular}
  \end{table}

  \begin{block}{اثبات ریاضی قانون پایستگی سیمپلکس در WGQA}
    مهار انفجار پرپلکسیتی از رقم نجومی \textbf{۴,۲۴۸,۶۸۹} به سطح متعارف با نگاشت بر روی سیمپلکس:
    \[
      \operatorname{scale}_h = \operatorname{Softmax}(w)_h \times H_{\text{kv}} \implies \sum_{h=1}^{H_{\text{kv}}} \operatorname{scale}_h = H_{\text{kv}}
    \]
    تضمین عدم‌تغییر انرژی سرهای KV و تثبیت کامل دمای لاجیت $T_{\text{eff}}$.
  \end{block}
\end{frame}

% --- Slide 7: ASVD KV-Cache Memory Scaling ---
\begin{frame}{شکاف دوم SOTA: فشرده‌سازی ۸ برابری بافر KV-Cache در ASVD}
  \begin{columns}
    \begin{column}{0.55\textwidth}
      \includegraphics[width=\textwidth]{kv_cache_latent_memory_scaling.png}
    \end{column}
    \begin{column}{0.45\textwidth}
      \begin{block}{هم‌ارزی استریمینگ و صرفه‌جویی حافظه}
        \begin{itemize}\footnotesize
          \item نگاشت کلیدها و مقادیر به بردار فشرده نهان مشترک $\mathbf{c}_t^{KV} \in \mathbb{R}^{d_c}$ با بعد $d_c \ll H d_h$.
          \item \textbf{روند مقیاس‌پذیری:} با افزایش طول توالی از ۲K به ۶۴K، مصرف بافر استاندارد از ۲۴ گیگابایت فراتر می‌رود، اما در ساختار ASVD در زیر ۳ گیگابایت باقی می‌ماند.
          \item بازسازی برخط در رجیسترهای پرسرعت بدون افت پرپلکسیتی.
        \end{itemize}
      \end{block}
    \end{column}
  \end{columns}
\end{frame}

% --- Slide 8: 4-Stage Surgical Injection Pipeline ---
\begin{frame}{پروتکل مهندسی تزریق جراحی ۴ مرحله‌ای}
  \begin{center}
    \begin{tikzpicture}[
      node distance=0.55cm,
      stage/.style={rectangle, draw=UTBlue, fill=LightBlue, rounded corners, minimum width=12cm, minimum height=0.85cm, font=\small},
      ar/.style={-Stealth, thick, draw=UTBlue}
    ]
      \node[stage] (s1) {\textbf{مرحله اول (Drop-in Surgical Injection):} جایگزینی درجا با حفظ کامل ابعاد $(B, T, D)$، ماسک علّی و RoPE/ALiBi};
      \node[stage, below=of s1] (s2) {\textbf{مرحله دوم (Parametric Alignment):} تراز طیفی مقادیر منفرد ($\tau \ge 0.90$) و مقداردهی PiSSA/DoRA};
      \node[stage, below=of s2] (s3) {\textbf{مرحله سوم (Fast Gradient Recovery):} بازیافت گرادیان در ۵۰ الی ۲۵۰ گام با بسته تقطیر ۶ ستونی در فضای پوچ هسین};
      \node[stage, below=of s3] (s4) {\textbf{مرحله چهارم (Hardware Profiling):} سنکرون‌سازی با رویدادهای CUDA، ۲۵ تکرار گرم‌کردن و استخراج توان عملیاتی};

      \draw[ar] (s1) -- (s2);
      \draw[ar] (s2) -- (s3);
      \draw[ar] (s3) -- (s4);
    \end{tikzpicture}
  \end{center}
  \vspace{-0.2cm}
  \begin{itemize}\footnotesize
    \item \textbf{ایزولاسیون کامل:} لایه‌های FFN، نرمال‌سازی‌ها و تعبیه‌ها فریز باقی می‌مانند.
    \item \textbf{برابری بازنمایی در گام صفر (Step-0 Parity):} جلوگیری از شوک ناگهانی به شبکه عصبی با انشعاب کم‌مرتبه.
  \end{itemize}
\end{frame}

% --- Slide 8b: 5-Stage Compression Pipeline Ablation Analysis ---
\begin{frame}{تحلیل ابلیشن ۵ مرحله‌ای خط لوله فشرده‌سازی روی DistilBERT (MNLI, $r=64$)}
  \begin{table}
    \centering
    \small
    \begin{tabular}{lccc}
      \toprule
      \textbf{پیکربندی / مرحله خط لوله} & \textbf{دقت (\%)} & \textbf{پیک حافظه (MB)} & \textbf{توان عملیاتی (tok/s)} \\
      \midrule
      ۱. فاکتورسازی ساده SVD (پایه کالیبره‌نشده) & 34.10\% & \textbf{163.0} & 51,244 \\
      ۲. + سفیدسازی چولسکی فعال‌سازی‌ها & 68.60\% & 178.5 & 57,900 \\
      ۳. + تخصیص پویا بر مبنای آنتروپی طیفی & 70.80\% & 172.4 & 58,400 \\
      ۴. + کِرنل فیوزشده خطی تریتون & 71.50\% & 168.2 & \textbf{61,200} \\
      ۵. + بازیافت خطای مانده ResSVD (کامل) & \textbf{72.20\%} & 170.2 & 59,200 \\
      \bottomrule
    \end{tabular}
  \end{table}

  \begin{block}{تحلیل مهندسی ابلیشن فشرده‌سازی}
    \begin{itemize}\footnotesize
      \item \textbf{جهش اولیه (+34.50\%):} سفیدسازی چولسکی فعال‌سازی‌ها از انحراف شدید کوواریانس جلوگیری کرده و دقت را بیش از دو برابر می‌کند ($34.10\% \to 68.60\%$).
      \item \textbf{مهار خطای انباشتی:} ترکیب تخصیص طیفی، کرنل فیوزشده و بازیابی مانده ResSVD دقت را به \textbf{72.20\%} ارتقا داده و بیشینه توان را تثبیت می‌نماید.
    \end{itemize}
  \end{block}
\end{frame}

% --- Slide 9: 6-Pillar Knowledge Distillation ---
\begin{frame}{بسته ۶ ستونی تقطیر دانش (6-Pillar KD Suite)}
  \begin{columns}
    \begin{column}{0.55\textwidth}
      \begin{block}{فرمول‌بندی تابع هدف مرکب تقطیر}
        \[
          \mathcal{L}_{\text{KD}}(\theta) = \sum_{i=1}^6 \lambda_i \mathcal{L}_i(\theta), \quad \sum_{i=1}^6 \lambda_i = 1, \quad \lambda_i > 0
        \]
        \begin{enumerate}\footnotesize
          \item \textbf{ستون اول:} تقطیر تفکیک‌شده لاجیت‌ها (DKD) و تطابق خطی CKA
          \item \textbf{ستون دوم:} هوموتوپی پیوسته توجه و تزریق تدریجی (MOHAWK)
          \item \textbf{ستون سوم:} آداپتورهای بهینه‌ساز زیرفضا (PiSSA و DoRA)
          \item \textbf{ستون چهارم:} ترازسازی ترجیحات بدون پاداش (SimPO / ORPO)
          \item \textbf{ستون پنجم:} پایدارسازهای عددی و سینک‌های توجه یادگیرنده
          \item \textbf{ستون ششم:} لایه انطباق در زمان آزمون (TTT Layer)
        \end{enumerate}
      \end{block}
    \end{column}

    \begin{column}{0.45\textwidth}
      \begin{exampleblock}{قضیه همگرایی انقباضی خطی}
        تحت فرض تحدب قوی $\mu = \sum \lambda_i \mu_i$ و صافی لیپ‌شیتز $L = \sum \lambda_i L_i$:
        \[
          G(\theta) = \theta - \gamma \nabla \mathcal{L}_{\text{KD}}(\theta)
        \]
        یک نگاشت انقباضی اکید با ضریب زیر است:
        \[
          \kappa^* = \sqrt{1 - \frac{\mu^2}{L^2}} < 1
        \]
        \[
          \|\theta_t - \theta^*\|_2 \le (\kappa^*)^t \|\theta_0 - \theta^*\|_2
        \]
        تضمین همگرایی هندسی پارامترها به وضعیت معلم.
      \end{exampleblock}
    \end{column}
  \end{columns}
\end{frame}

% --- Slide 10: Empirical Convergence & Distillation Results ---
\begin{frame}{اثربخشی تجربی بسته ۶ ستونی در بازیافت دقت}
  \begin{columns}
    \begin{column}{0.5\textwidth}
      \includegraphics[width=\textwidth]{chart_training_loss_convergence.png}
    \end{column}
    \begin{column}{0.5\textwidth}
      \begin{table}
        \centering
        \tiny
        \begin{tabular}{llccc}
          \toprule
          \textbf{ماژول} & \textbf{رژیم آموزشی} & \textbf{دقت} & \textbf{انتروپی} & \textbf{توان (tok/s)} \\
          \midrule
          \multirow{3}{*}{\textbf{Cosformer}} & بدون آموزش (Zero-Shot) & 0.5519 & 2.9356 & 488.9 \\
          & انتروپی متقاطع خام (CE) & 0.2286 & 4.4928 & 566.9 \\
          & \textbf{بسته ۶ ستونی (6-Pillar)} & \textbf{0.9184} & \textbf{0.5125} & \textbf{1029.2} \\
          \midrule
          \multirow{3}{*}{\textbf{Chunk GLA}} & بدون آموزش (Zero-Shot) & 0.0994 & 5.8951 & 1029.2 \\
          & انتروپی متقاطع خام (CE) & 0.1578 & 4.8950 & 551.6 \\
          & \textbf{بسته ۶ ستونی (6-Pillar)} & \textbf{0.7871} & \textbf{1.2447} & \textbf{1014.4} \\
          \midrule
          \multirow{3}{*}{\textbf{Mamba Block}} & بدون آموزش (Zero-Shot) & 0.8416 & 1.0345 & 1144.2 \\
          & انتروپی متقاطع خام (CE) & 0.2223 & 4.5791 & 643.4 \\
          & \textbf{بسته ۶ ستونی (6-Pillar)} & \textbf{0.9827} & \textbf{0.1124} & \textbf{1105.7} \\
          \bottomrule
        \end{tabular}
      \end{table}
      \begin{block}{\small یافته تجربی کلیدی}
        انتروپی متقاطع خام به دلیل انحراف شدید بازنمایی دچار واگرایی می‌شود، در حالی که بسته ۶ ستونی با هدایت زیرفضای هسین، بازیابی بیش از ۹۰٪ دقت معلم را محقق می‌سازد.
      \end{block}
    \end{column}
  \end{columns}
\end{frame}

% --- Slide 11: Large-Scale Benchmark & Taxonomy ---
\begin{frame}{گستره بنچ‌مارک تجربی: ۲۰۳ ماژول در ۶ خانواده معماری}
  \begin{columns}
    \begin{column}{0.5\textwidth}
      \includegraphics[width=\textwidth]{chart_radar_families.png}
    \end{column}
    \begin{column}{0.5\textwidth}
      \begin{block}{تقسیم‌بندی ماژول‌های مورد ارزیابی}
        \begin{itemize}\footnotesize
          \setlength{\itemsep}{1pt}
          \item \textbf{۱۳۱ مکانیزم توجه متنوع:}
            \begin{itemize}\scriptsize
              \item توجه خطی و بازگشتی (Linear, Cosformer, GLA, DeltaNet, RWKV)
              \item ساختار خلوت و پنجره‌ای (Longformer, BigBird, NSA/DSA)
              \item مدل‌های فضای حالت (Mamba, S4, H3)
            \end{itemize}
          \item \textbf{۷۲ موتور فشرده‌سازی و تجزیه کم‌مرتبه:}
            \begin{itemize}\scriptsize
              \item روش‌های SVD پیشرفته (ASVD, SVD-LLM v2, AASVD, PiSSA, DoRA)
              \item کوانتایزاسیون ماتریس‌های منفرد (FP8, INT4, NF4)
              \item هرس طیفی و خبرگان (REAP-MoE, Structural Pruning)
            \end{itemize}
          \item \textbf{۴ مدل پایه ناهمگون:}
            \begin{itemize}\scriptsize
              \item انکودرها: DistilBERT, RoBERTa
              \item دیکودرهای علّی: GPT-2, Qwen-0.5B
            \end{itemize}
        \end{itemize}
      \end{block}
    \end{column}
  \end{columns}
\end{frame}

% --- Slide 12: Top-15 Accuracy Ranking ---
\begin{frame}{رتبه‌بندی ۱۵ ماژول پیشتاز در دقت زبانی}
  \begin{columns}
    \begin{column}{0.65\textwidth}
      \includegraphics[width=\textwidth]{chart_top15_accuracy.png}
    \end{column}
    \begin{column}{0.35\textwidth}
      \begin{block}{نکات تحلیلی}
        \begin{itemize}\footnotesize
          \item برتری قاطع معماری‌های ترکیب خطی با تراز طیفی (نظیر \lr{Cosformer+AASVD} و \lr{Mamba Block}).
          \item حفظ بیش از ۹۵٪ دقت مدل متراکم اولیه در کسری از حافظه مصرفی.
          \item برتری پایدار مدل‌های مجهز به رژیم ۶ ستونی در قیاس با رقبا.
        \end{itemize}
      \end{block}
    \end{column}
  \end{columns}
\end{frame}

% --- Slide 13: Pareto Frontier (Accuracy vs VRAM) ---
\begin{frame}{جبهه بهینگی پارتو: دقت در برابر اوج مصرف حافظه (VRAM)}
  \begin{columns}
    \begin{column}{0.62\textwidth}
      \includegraphics[width=\textwidth]{pareto_accuracy_vs_vram.png}
    \end{column}
    \begin{column}{0.38\textwidth}
      \begin{exampleblock}{نقاط بهینه جبهه پارتو}
        \begin{itemize}\footnotesize
          \item \textbf{رژیم بیشینه فشرده‌سازی:} ماژول‌های \lr{ASVD} با فشرده‌سازی ۸ برابری بافر، مصرف حافظه را به زیر ۶۰ مگابایت می‌رسانند.
          \item \textbf{رژیم توازن طلایی:} مدل \lr{Cosformer+AASVD} با ردپای ۶۷ مگابایت و دقت ۹۴.۳٪ در راس جبهه پارتو جای می‌گیرد.
          \item \textbf{رژیم بیشینه دقت:} \lr{Mamba} و \lr{GDN-2} با دقت بالای ۹۸٪ و کارایی حافظه عالی.
        \end{itemize}
      \end{exampleblock}
    \end{column}
  \end{columns}
\end{frame}

% --- Slide 14: Pareto Frontier (Accuracy vs Throughput) ---
\begin{frame}{جبهه بهینگی پارتو: دقت در برابر توان پردازشی (Throughput)}
  \begin{columns}
    \begin{column}{0.62\textwidth}
      \includegraphics[width=\textwidth]{pareto_accuracy_vs_throughput.png}
    \end{column}
    \begin{column}{0.38\textwidth}
      \begin{block}{تحلیل توان عملیاتی توکن بر ثانیه}
        \begin{itemize}\footnotesize
          \item مدل متراکم اولیه به دلیل سربار $\mathcal{O}(N^2)$ به نرخ ۵۰۰ توکن بر ثانیه محدود است.
          \item ماژول‌های خطی بهینه‌شده به همراه ادغام کرنل، توان را به بیش از \textbf{۱۱۰۰ توکن بر ثانیه} ارتقا می‌دهند ($>2.2\times$).
          \item مدل‌های Mamba و DeltaNet بالاترین توان پردازشی را با دقت بالای ۹۵٪ ثبت کرده‌اند.
        \end{itemize}
      \end{block}
    \end{column}
  \end{columns}
\end{frame}

% --- Slide 15: Sequence Length Scaling ---
\begin{frame}{رفتار مقیاس‌پذیری با افزایش طول توالی ($N = 128 \to 8192$)}
  \begin{columns}
    \begin{column}{0.6\textwidth}
      \includegraphics[width=\textwidth]{sequence_length_scaling.png}
    \end{column}
    \begin{column}{0.4\textwidth}
      \begin{alertblock}{شکستن دیوار محاسباتی}
        \begin{itemize}\footnotesize
          \item در مکانیزم توجه استاندارد \lr{Dense Softmax}، تاخیر با شیب درجه دوم $\mathcal{O}(N^2)$ افزایش یافته و در $N=8192$ دچار انفجار تاخیر یا OOM می‌شود.
          \item مکانیزم‌های خطی و ساختاریافته رفتار خطی اکید $\mathcal{O}(N)$ را حفظ نموده و تاخیر را در توالی‌های بسیار بلند زیر ۱۰ میلی‌ثانیه نگه می‌دارند.
        \end{itemize}
      \end{alertblock}
    \end{column}
  \end{columns}
\end{frame}

% --- Slide 16: Hardware Acceleration with 26 Triton Kernels ---
\begin{frame}{بهینه‌سازی سطح پایین: ۲۶ کرنل سفارشی Triton JIT روی RTX 4070}
  \begin{columns}
    \begin{column}{0.48\textwidth}
      \begin{block}{مشخصات پلتفرم تست سخت‌افزاری}
        \begin{itemize}\footnotesize
          \item پردازنده گرافیکی: \textbf{NVIDIA GeForce RTX 4070 Laptop}
          \item معماری سخت‌افزار: \textbf{Ada Lovelace (Compute Capability 8.9)}
          \item حافظه فیزیکی: \textbf{8188 MB GDDR6}
          \item کامپایلر: \textbf{OpenAI Triton 3.0.0 JIT}
        \end{itemize}
      \end{block}

      \begin{alertblock}{پشته گسیل سه‌سطحی (3-Level Dispatch)}
        \begin{enumerate}\footnotesize
          \item \textbf{Native Triton:} اجرای کاشی‌بندی‌شده در SRAM.
          \item \textbf{CUDA C++:} ۱۸ عملگر سفارشی با warp intrinsics.
          \item \textbf{PyTorch Fallback:} پشتیبان عددی تضمین‌کننده پایداری.
        \end{enumerate}
      \end{alertblock}
    \end{column}

    \begin{column}{0.52\textwidth}
      \includegraphics[width=\textwidth]{triton_speedup_leaderboard.png}
    \end{column}
  \end{columns}
\end{frame}

% --- Slide 17: Three-Level Latency Comparison ---
\begin{frame}{مقایسه تاخیر در پشته گسیل سه‌سطحی (PyTorch vs CUDA vs Triton)}
  \begin{columns}
    \begin{column}{0.6\textwidth}
      \includegraphics[width=\textwidth]{three_level_latency_comparison.png}
    \end{column}
    \begin{column}{0.4\textwidth}
      \begin{exampleblock}{دستاوردهای کرنل‌های تریتون}
        \begin{itemize}\footnotesize
          \item شتاب تا \textbf{۳.۸۱ برابر} در ماژول \lr{Chunk GLA}.
          \item شتاب \textbf{۳.۴۸ برابر} در محاسبات پروجکشن کم‌مرتبه \lr{SVD}.
          \item حذف کامل گلوگاه حافظه با ادغام عملگرهای نرمال‌سازی و رزیدوال (\lr{RMSNorm + Residual}).
          \item نزدیک شدن شدت محاسباتی به سقف سخت‌افزاری پردازنده RTX 4070.
        \end{itemize}
      \end{exampleblock}
    \end{column}
  \end{columns}
\end{frame}

% --- Slide 18: Contributions & Summary ---
\begin{frame}{جمع‌بندی دستاوردها و مشارکت‌های پایان‌نامه}
  \begin{enumerate}\setlength{\itemsep}{0.22cm}\footnotesize
    \item \textbf{بسته‌شدن قطعی ۵ شکاف علمی در مرز دانش (SOTA):} اثبات تحلیلی بهینگی برش چولسکی در SVD-LLM v2، استریمینگ بافر KV در ASVD، کران خطای مانتیس FP8، هم‌ترازی DSA، و پایستگی سیمپلکس در WGQA.
    \item \textbf{جامع‌ترین پایگاه ارزیابی مقایسه‌ای هم‌شرط:} سنجش دقیق ۲۰۳ ماژول (۱۳۱ توجه و ۷۲ فشرده‌ساز) با سنکرون‌سازی کامل CUDA روی ۴ مدل انکودر و دیکودر.
    \item \textbf{طراحی پروتکل مهندسی تزریق جراحی ۴ مرحله‌ای:} چارچوب سیستم‌محور برای جایگزینی درجا با برابری خروجی در گام صفر.
    \item \textbf{بسته ۶ ستونی تقطیر دانش با تضمین انقباض خطی:} مهار واگرایی بازنمایی و بازیابی بیش از ۹۰٪ دقت اولیه در کمتر از ۲۵۰ گام بهینه‌سازی.
    \item \textbf{شتاب‌دهی سخت‌افزاری سطح پایین:} توسعه و کامپایل ۲۶ کرنل سخت‌افزاری Triton JIT با شتاب تا ۳.۸ برابر روی \lr{RTX 4070 SM 8.9}.
  \end{enumerate}
\end{frame}

% --- Slide 19: Q&A Slide ---
\begin{frame}[plain]
  \begin{center}
    \vspace{0.8cm}
    \includegraphics[height=1.6cm]{logo.png}\\[0.5cm]
    {\huge \color{UTBlue} \textbf{با سپاس از توجه اساتید و داوران گرامی}}\\[0.8cm]
    {\Large \textbf{پرسش و پاسخ}}\\[1.2cm]
    {\small محمد طاها مجلسی --- دانشکده مهندسی برق و کامپیوتر دانشگاه تهران}\\[0.2cm]
    {\footnotesize \lr{Email: taha.majlesi@ut.ac.ir}}
  \end{center}
\end{frame}

\end{document}
